% T20 source: discussion.tex lines 58--119
\subsection{Relation to Prior Work}

The human-resource ceiling is a structural analogue of the logic of Amdahl's
law: the sequential component here is not a segment of a computational program
but the measured involvement of the developer \cite{amdahl1967validity}. An
interior
minimum over configured $P$ is consistent with the Universal Scalability Law in
the limited sense that increasing contention and coordination can turn
saturation into negative returns \cite{gunther2008general}. The numerical form
of $C(P)$ for agents nevertheless remains a hypothesis of the model, not a
consequence of these classical results.

The gap between $B_4$ and $T^*$ connects the formulation to scheduling theory
and the RCPSP: the original task-dependency graph and aggregate load are
insufficient when a single service resource creates additional chains induced
by resource constraints \cite{graham1966bounds,brucker1999resource}.
Resource-order arcs and a shared server are standard mechanisms in scheduling
theory \cite{pinedo2008scheduling,hall2000parallel}. Loading and unloading models
already include service before and after processing and may retain a machine
until unloading \cite{elidrissi2023loading}; the reentrant model routes work back
to the same primary machine after a remote server but releases that machine
during service \cite{chakhlevitch2009reentrant}. Accordingly, no novelty is
claimed for DAGs, a shared resource, repeated service, or a resource-augmented
graph as such.

A related concept of a finite number of robots that can be controlled
effectively appears in models of the number of robots per operator, in which
service intervals must fit within the autonomous-work windows of the other
systems \cite{olsen2004fanout,crandall2005validating}. The experiment by Cummings
and Mitchell further shows that queueing, delays in regaining situational
awareness, and retasking reduce the operator's predicted throughput, but its
UAV-simulation results cannot be transferred numerically to software development
\cite{cummings2008predicting}. Among coding-agent systems, MAGIS implements an
automated ``development---quality control---revision'' cycle, while CAID uses
concurrent worktrees with a dependency plan, merging, and testing
\cite{tao2024magis,geng2026caid}. Their quality and success metrics do not
demonstrate a reduction in makespan; in CAID's main comparisons, elapsed time
and cost increased together with the quality metric.

The present model combines known constructs in a mapping to the software
process: a task- and mode-dependent coefficient $k$, separate $a_i$ and $h_i$,
repeated human-attention phases, local retention of an agent slot, integration
and attention-switching overhead, and the makespan of an accepted result. The
main reviewed corpus, assembled under a search protocol with a cutoff date of
2026-08-09, contained no joint empirical calibration of these elements on
software tasks; this is not a universal claim about the entire literature. The
present article likewise does not close this empirical gap: the computational
results clarify the mechanism of the process constraint in synthetic scenarios
but do not replace measurements of specific tools and teams.

\subsection{Using the Model}

First, tasks and operating modes should be estimated on a common scale, and the
task-DAG and phase profiles should be constructed; $B_4$ can then be computed as
a rapid bottleneck diagnostic. For several competing configurations, an exact
or heuristic schedule is constructed, and whether a certified $T^*$ or only
$T(\pi)$ was obtained is recorded explicitly. Finally, sensitivity to $k$, $h$,
$C$, and $\gamma$ is evaluated. The model is more useful for comparing variants
of work organization than for promising a precise elapsed time without local
calibration. The active branch of $B_4$ characterizes the constraint that
determines the lower bound, but does not by itself prove that $T^*$ is governed
by the same branch: such a conclusion requires either
attainability of the bound or analysis of the resource-augmented schedule.
Similarly, an autonomous agent phase means only that the developer is not
involved during that interval; a task containing such a phase remains delegated
if task specification or acceptance requires human involvement.
